
\documentclass{article}
\usepackage{colortbl}
\usepackage{makecell}
\usepackage{multirow}
\usepackage{supertabular}

\begin{document}

\newcounter{utterance}

\centering \large Interaction Transcript for game `dond', experiment `coop\_en', episode 1 with reflexion{-}Qwen3.5{-}27B.
\vspace{24pt}

{ \footnotesize  \setcounter{utterance}{1}
\setlength{\tabcolsep}{0pt}
\begin{supertabular}{c@{$\;$}|p{.15\linewidth}@{}p{.15\linewidth}p{.15\linewidth}p{.15\linewidth}p{.15\linewidth}p{.15\linewidth}}
    \# & $\;$A & \multicolumn{4}{c}{Game Master} & $\;\:$B\\
    \hline

    \theutterance \stepcounter{utterance}  
    & & \multicolumn{4}{p{0.6\linewidth}}{
        \cellcolor[rgb]{0.9,0.9,0.9}{
            \makecell[{{p{\linewidth}}}]{
                \texttt{\tiny{[P1$\langle$GM]}}
                \texttt{You have attempted this same task before and did not succeed. These are the plans you wrote after those attempts.} \\
\texttt{Trial \#1: New plan: I failed because I submitted a secret proposal prematurely without explicitly confirming the other player's specific allocation in their own proposal format, leading to a potential mismatch or lack of mutual agreement on the final split. Although we verbally agreed on the distribution, the game requires both parties to submit complementary secret proposals to succeed. Next time, after reaching a verbal consensus, I will explicitly ask the other player to submit their secret proposal first or confirm their exact list before I submit mine. Specifically, I will wait for the other player to send a message containing "{[}Proposal: ...{]}" or explicitly state "I am submitting my proposal now" before I send my own. If they do not submit a proposal after my verbal confirmation, I will reiterate the agreed split and ask them to confirm their specific list in the proposal format to ensure our proposals are perfectly complementary (e.g., ensuring they propose taking the cup and flower while I propose taking the remote, box, and frame). I will not submit my secret proposal until I have high confidence that the other player has also committed to the complementary set.} \\
\\ 
\\ 
\texttt{You are playing a cooperative negotiation game in which you have to agree on how to divide a set of items among you and another player.} \\
\\ 
\texttt{Rules:} \\
\texttt{(a) You and the other player are given a set of items. Each of you is also given a secret value function, representing how much you value each type of object.} \\
\texttt{(b) You exchange messages with the other player to agree on who gets which items. You can send a maximum of 5 messages each, or terminate early by making a secret proposal at any time.} \\
\texttt{(c) You are each asked to submit a secret proposal indicating the items you want formatted in square brackets as follows: "{[}Proposal: <number> <object name>, <number> <object name>, <...>{]}"} \\
\texttt{(d) If your proposals are complementary, i.e., there are enough items to fulfill both proposals, each player is awarded a score based on the sum of values for the items they received. Otherwise, both of you get zero points.} \\
\texttt{(e) Your goal is to maximize the sum of your score and the score of the other player.} \\
\\ 
\texttt{Let us start.} \\
\\ 
\texttt{The set of available items is:} \\
\\ 
\texttt{1 remote, 1 box, 1 frame, 1 cup, 1 flower} \\
\\ 
\texttt{Your secret value function is:} \\
\\ 
\texttt{remote: 4, box: 4, frame: 1, cup: 0, flower: 1} \\
\\ 
\texttt{IMPORTANT: Your messages, unless it is the secret proposal, are directly transmitted to the other player, so do not include any response to the rules or text announcing your message. To make a secret proposal, use the indicated format. Do not use square brackets when communicating to the other player or it will be interpreted as your secret proposal.} \\
            }
        }
    }
    & \\ \\

    \theutterance \stepcounter{utterance}  
    & \multicolumn{4}{p{0.6\linewidth}}{
        \cellcolor[rgb]{0.9,0.9,0.9}{
            \makecell[{{p{\linewidth}}}]{
                \texttt{\tiny{[P1$\rangle$GM]}}
                \texttt{Hello! Let's work together to divide these items efficiently. Based on my values, I really need the remote and the box (I value them highly), while I don't care much about the cup, frame, or flower.} \\
\\ 
\texttt{Could you share your secret values or tell me which items you are most interested in? That way, we can make sure we both get the items we value most and avoid any conflicts.} \\
            }
        }
    }
    & & \\ \\

    \theutterance \stepcounter{utterance}  
    & & & \multicolumn{2}{p{0.3\linewidth}}{
        \cellcolor[rgb]{0.9,0.9,0.9}{
            \makecell[{{p{\linewidth}}}]{
                \texttt{\tiny{[GM$|$GM]}}
                \texttt{continue} \\
            }
        }
    }
    & & \\ \\

    \theutterance \stepcounter{utterance}  
    & & & \multicolumn{2}{p{0.3\linewidth}}{
        \cellcolor[rgb]{0.9,0.9,0.9}{
            \makecell[{{p{\linewidth}}}]{
                \texttt{\tiny{[GM$|$GM]}}
                \texttt{Hello! Let's work together to divide these items efficiently. Based on my values, I really need the remote and the box (I value them highly), while I don't care much about the cup, frame, or flower.} \\
\\ 
\texttt{Could you share your secret values or tell me which items you are most interested in? That way, we can make sure we both get the items we value most and avoid any conflicts.} \\
            }
        }
    }
    & & \\ \\

    \theutterance \stepcounter{utterance}  
    & & \multicolumn{4}{p{0.6\linewidth}}{
        \cellcolor[rgb]{0.9,0.9,0.9}{
            \makecell[{{p{\linewidth}}}]{
                \texttt{\tiny{[GM$\rangle$P2]}}
                \texttt{You have attempted this same task before and did not succeed. These are the plans you wrote after those attempts.} \\
\texttt{Trial \#1: New plan: I failed because I ignored my own secret value function and the negotiation logic I established. In my message, I correctly identified that I value the cup (6) and the other player values the remote (4), yet my final proposal claimed the remote and box (which I value at 0) while giving up the cup. This resulted in a non{-}complementary or low{-}value outcome. Next time, I will strictly align my proposal with my highest{-}value items (cup) and the other player's highest{-}value items (remote, box) as discussed. I will propose: "{[}Proposal: 1 cup, 1 flower{]}" (or frame), ensuring the other player gets the remote and box, thereby maximizing the combined score and ensuring the proposals are complementary.} \\
\\ 
\\ 
\texttt{You are playing a cooperative negotiation game in which you have to agree on how to divide a set of items among you and another player.} \\
\\ 
\texttt{Rules:} \\
\texttt{(a) You and the other player are given a set of items. Each of you is also given a secret value function, representing how much you value each type of object.} \\
\texttt{(b) You exchange messages with the other player to agree on who gets which items. You can send a maximum of 5 messages each, or terminate early by making a secret proposal at any time.} \\
\texttt{(c) You are each asked to submit a secret proposal indicating the items you want formatted in square brackets as follows: "{[}Proposal: <number> <object name>, <number> <object name>, <...>{]}"} \\
\texttt{(d) If your proposals are complementary, i.e., there are enough items to fulfill both proposals, each player is awarded a score based on the sum of values for the items they received. Otherwise, both of you get zero points.} \\
\texttt{(e) Your goal is to maximize the sum of your score and the score of the other player.} \\
\\ 
\texttt{Let us start.} \\
\\ 
\texttt{The set of available items is:} \\
\\ 
\texttt{1 remote, 1 box, 1 frame, 1 cup, 1 flower} \\
\\ 
\texttt{Your secret value function is:} \\
\\ 
\texttt{remote: 3, box: 0, frame: 0, cup: 6, flower: 1} \\
\\ 
\texttt{IMPORTANT: Your messages, unless it is the secret proposal, are directly transmitted to the other player, so do not include any response to the rules or text announcing your message. To make a secret proposal, use the indicated format. Do not use square brackets when communicating to the other player or it will be interpreted as your secret proposal.} \\
\\ 
\\ 
\texttt{Hello! Let's work together to divide these items efficiently. Based on my values, I really need the remote and the box (I value them highly), while I don't care much about the cup, frame, or flower.} \\
\\ 
\texttt{Could you share your secret values or tell me which items you are most interested in? That way, we can make sure we both get the items we value most and avoid any conflicts.} \\
            }
        }
    }
    & \\ \\

    \theutterance \stepcounter{utterance}  
    & & & \multicolumn{4}{p{0.6\linewidth}}{
        \cellcolor[rgb]{0.9,0.9,0.9}{
            \makecell[{{p{\linewidth}}}]{
                \texttt{\tiny{[GM$\langle$P2]}}
                \texttt{Hello! Thanks for sharing your preferences. To help us find the best split, here are my values: I value the **cup** the most (6 points), followed by the **remote** (3 points), and the **flower** (1 point). I don't value the **box** or the **frame** at all (0 points).} \\
\\ 
\texttt{Since you really need the remote and the box, and I really need the cup, it seems like we can easily avoid conflict. If you take the **remote** and the **box**, and I take the **cup**, we both get our highest priorities. We can then split the **frame** and **flower** however you prefer, or I can take the flower since I value it slightly.} \\
\\ 
\texttt{Does a deal where you get the remote and box, and I get the cup (and maybe the flower), sound good to you?} \\
            }
        }
    }
     \\ \\

    \theutterance \stepcounter{utterance}  
    & & & \multicolumn{2}{p{0.3\linewidth}}{
        \cellcolor[rgb]{0.9,0.9,0.9}{
            \makecell[{{p{\linewidth}}}]{
                \texttt{\tiny{[GM$|$GM]}}
                \texttt{continue} \\
            }
        }
    }
    & & \\ \\

    \theutterance \stepcounter{utterance}  
    & & & \multicolumn{2}{p{0.3\linewidth}}{
        \cellcolor[rgb]{0.9,0.9,0.9}{
            \makecell[{{p{\linewidth}}}]{
                \texttt{\tiny{[GM$|$GM]}}
                \texttt{Hello! Thanks for sharing your preferences. To help us find the best split, here are my values: I value the **cup** the most (6 points), followed by the **remote** (3 points), and the **flower** (1 point). I don't value the **box** or the **frame** at all (0 points).} \\
\\ 
\texttt{Since you really need the remote and the box, and I really need the cup, it seems like we can easily avoid conflict. If you take the **remote** and the **box**, and I take the **cup**, we both get our highest priorities. We can then split the **frame** and **flower** however you prefer, or I can take the flower since I value it slightly.} \\
\\ 
\texttt{Does a deal where you get the remote and box, and I get the cup (and maybe the flower), sound good to you?} \\
            }
        }
    }
    & & \\ \\

    \theutterance \stepcounter{utterance}  
    & & \multicolumn{4}{p{0.6\linewidth}}{
        \cellcolor[rgb]{0.9,0.9,0.9}{
            \makecell[{{p{\linewidth}}}]{
                \texttt{\tiny{[P1$\langle$GM]}}
                \texttt{Hello! Thanks for sharing your preferences. To help us find the best split, here are my values: I value the **cup** the most (6 points), followed by the **remote** (3 points), and the **flower** (1 point). I don't value the **box** or the **frame** at all (0 points).} \\
\\ 
\texttt{Since you really need the remote and the box, and I really need the cup, it seems like we can easily avoid conflict. If you take the **remote** and the **box**, and I take the **cup**, we both get our highest priorities. We can then split the **frame** and **flower** however you prefer, or I can take the flower since I value it slightly.} \\
\\ 
\texttt{Does a deal where you get the remote and box, and I get the cup (and maybe the flower), sound good to you?} \\
            }
        }
    }
    & \\ \\

    \theutterance \stepcounter{utterance}  
    & \multicolumn{4}{p{0.6\linewidth}}{
        \cellcolor[rgb]{0.9,0.9,0.9}{
            \makecell[{{p{\linewidth}}}]{
                \texttt{\tiny{[P1$\rangle$GM]}}
                \texttt{That sounds like a perfect plan! Since I value the remote and box the most (4 points each) and you value the cup the most (6 points), giving those to us respectively maximizes our individual scores.} \\
\\ 
\texttt{Regarding the remaining items:} \\
\texttt{{-} I value the frame (1 point) and flower (1 point).} \\
\texttt{{-} You value the flower (1 point) and frame (0 points).} \\
\\ 
\texttt{To maximize our **combined** score, it makes the most sense for me to take the **frame** (since I get 1 point and you get 0) and for you to take the **flower** (since you get 1 point and I get 1 point, but wait—if I take the flower I get 1, if you take it you get 1. It's a tie in total value, but let's look at the split).} \\
\\ 
\texttt{Actually, let's re{-}evaluate the frame and flower to be absolutely sure we are maximizing the *sum*:} \\
\texttt{{-} **Option A**: I take Frame (1) + Flower (1) = 2 points for me. You take nothing extra. Total extra = 2.} \\
\texttt{{-} **Option B**: I take Frame (1). You take Flower (1). Total extra = 2.} \\
\texttt{{-} **Option C**: I take Flower (1). You take Frame (0). Total extra = 1.} \\
\\ 
\texttt{So, Option A and Option B both yield the same total sum (2 points). However, since you mentioned you value the flower slightly, let's give you the **flower** and I will take the **frame**. This way, we both get a small bonus, and the total sum is maximized.} \\
\\ 
\texttt{**Proposed Final Split:**} \\
\texttt{{-} **You**: Cup, Flower} \\
\texttt{{-} **Me**: Remote, Box, Frame} \\
\\ 
\texttt{Does this allocation work for you? If so, please confirm, and we can proceed to submit our secret proposals. To ensure we don't fail like last time, could you please send your secret proposal first using the format `{[}Proposal: ...{]}` so I can match it exactly with mine?} \\
            }
        }
    }
    & & \\ \\

    \theutterance \stepcounter{utterance}  
    & & & \multicolumn{2}{p{0.3\linewidth}}{
        \cellcolor[rgb]{0.9,0.9,0.9}{
            \makecell[{{p{\linewidth}}}]{
                \texttt{\tiny{[GM$|$GM]}}
                \texttt{proposal does not include number/name pairs} \\
            }
        }
    }
    & & \\ \\

    \theutterance \stepcounter{utterance}  
    & \multicolumn{4}{p{0.6\linewidth}}{
        \cellcolor[rgb]{0.95,0.95,0.95}{
            \makecell[{{p{\linewidth}}}]{
                \texttt{\tiny{[P2$\rightarrow$P1]}}
                \texttt{Post{-}episode Reflexion (not sent to GM):} \\
\texttt{New plan: The previous attempt was aborted because I included square brackets in my message to the other player ("{[}Proposal: ...{]}"), which the game rules explicitly forbid for non{-}secret proposals. The system interpreted this as an invalid move or a premature secret proposal attempt. Next time, I will strictly avoid using square brackets in any conversational message. Instead, I will verbally confirm the agreed split (e.g., "I will take Remote, Box, Frame; you take Cup, Flower") and then immediately submit my secret proposal using the correct format `{[}Proposal: 1 remote, 1 box, 1 frame{]}` without asking the other player to send their proposal first or using brackets in the negotiation text. I will ensure the proposal format is only used for the final submission step.} \\
            }
        }
    }
    & & \\ \\

    \theutterance \stepcounter{utterance}  
    & & & \multicolumn{4}{p{0.6\linewidth}}{
        \cellcolor[rgb]{0.95,0.95,0.95}{
            \makecell[{{p{\linewidth}}}]{
                \texttt{\tiny{[P1$\rightarrow$P2]}}
                \texttt{Post{-}episode Reflexion (not sent to GM):} \\
\texttt{New plan: The previous attempt was aborted because I violated the rule prohibiting square brackets in negotiation messages. I used square brackets to highlight item names (e.g., "{[}Proposal: ...{]}" or bolding with brackets) in my text to the other player, which the system interpreted as a premature secret proposal. Next time, I will strictly avoid using square brackets in any message sent to the other player. I will communicate my values and proposed split using plain text or standard punctuation (like parentheses or bolding without brackets) to ensure the message is treated as a negotiation step. Only when I am ready to finalize the deal will I use the exact format "{[}Proposal: <number> <object name>, ...{]}" as a standalone message.} \\
            }
        }
    }
     \\ \\

\end{supertabular}
}

\end{document}
